%% Generated from the running pgm-studio API (http://localhost:7894/api) --- do not edit by hand.
%% Source: GET /api/openapi/v1.json
% 231 operations over 174 paths, grouped by the area segment of the path.

\begin{longtable}{@{}L{14mm}L{47mm}L{79mm}@{}}
\toprule
\textbf{Method} & \textbf{Path} & \textbf{Summary} \\
\midrule
\endfirsthead
\multicolumn{3}{@{}l}{\footnotesize\itshape\color{ink60}Endpoint surface, continued}\\[2pt]
\toprule
\textbf{Method} & \textbf{Path} & \textbf{Summary} \\
\midrule
\endhead
\midrule
\multicolumn{3}{r@{}}{\footnotesize\color{ink60}continued}\\
\endfoot
\bottomrule
\endlastfoot
\multicolumn{3}{@{}l}{\textbf{\id{/map}} \textcolor{ink60}{\small(126)}}\\[1pt]
\midrule[0.2pt]
\id{POST} & \id{/\allowbreak{}api/\allowbreak{}map/\allowbreak{}from-documents} & load a whole map from the three documents it is made of: the plan it was drawn from, the layout that was built, and the intent it is played for. The way back in for a map authored elsewhere.\,\textcolor{ink30}{\ldots} \\
\id{POST} & \id{/\allowbreak{}api/\allowbreak{}map/\allowbreak{}import-folder} & import a local xml-less world (B8 ``open a folder''): resolve \textless{}root\textgreater{}/\textless{}slug\textgreater{}/region via MapsRoots (only configured roots --- no client path), create the map row, and scan into MariaDB.\,\textcolor{ink30}{\ldots} \\
\id{POST} & \id{/\allowbreak{}api/\allowbreak{}map/\allowbreak{}import-url} & B8 import-from-url (docs/tools/configure.md, the Import phase). Server-side: fetch a zipped Minecraft world from an allowlisted host, safely extract only region/.mca, create the map row, and scan it into MariaDB (reusing WorldFeatureWriter).\,\textcolor{ink30}{\ldots} \\
\id{GET} & \id{/\allowbreak{}api/\allowbreak{}map/\allowbreak{}\{slug\}} & the full map document (the xml\_data.json shape), reconstructed from the relational rows. \\
\id{GET} & \id{/\allowbreak{}api/\allowbreak{}map/\allowbreak{}\{slug\}/\allowbreak{}block-seat} & ?x=\&y=\&z= --- whether one block can hold a thing placed into it, read off the same vertical segments column-floor takes. A monument is the block a player puts a wool into, so the authoring step asks both halves at once: the block itself must\,\textcolor{ink30}{\ldots} \\
\id{GET} & \id{/\allowbreak{}api/\allowbreak{}map/\allowbreak{}\{slug\}/\allowbreak{}column} & One or more columns bedrock-to-sky, every block named, as text. THE WORKHORSE: every picture beside it is a projection, and this is what is actually at a coordinate --- which is why it is the read to reach for when a picture and a document disagree.\,\textcolor{ink30}{\ldots} \\
\id{GET} & \id{/\allowbreak{}api/\allowbreak{}map/\allowbreak{}\{slug\}/\allowbreak{}column-floor} & ?x=\&z=\&y= --- the terrain floor Y at a single column from the vertical segments: the top of the highest solid segment at or below the reference y (the floor a thing at y rests on), falling back to the segment top nearest y.\,\textcolor{ink30}{\ldots} \\
\id{GET} & \id{/\allowbreak{}api/\allowbreak{}map/\allowbreak{}\{slug\}/\allowbreak{}core-suggestions} & [?box=x0,y0,z0,x1,y1,z1] --- the cores the ingest scan proposed for this map, read straight out of core\_candidate. No world access and no scoring pass: a core's signature is unambiguous enough that the gather already knew the casing, so the\,\textcolor{ink30}{\ldots} \\
\id{GET} & \id{/\allowbreak{}api/\allowbreak{}map/\allowbreak{}\{slug\}/\allowbreak{}coverage} & where the ground is lived on and where it is dead: the traffic corridors between every pair of waypoints, each waypoint's own ring, the prop-decorated fringe, and the dead patches named with coordinates. ?format=png answers the same grid as the coverage picture.\,\textcolor{ink30}{\ldots} \\
\id{GET} & \id{/\allowbreak{}api/\allowbreak{}map/\allowbreak{}\{slug\}/\allowbreak{}destroyable-suggestions} & [?box=x0,y0,z0,x1,y1,z1] --- the destroyables the ingest scan proposed for this map, read straight out of destroyable\_candidate. box is optional and narrows the list to masses intersecting it, which is how the configure step scopes a detect to the area the author drew.\,\textcolor{ink30}{\ldots} \\
\id{GET} & \id{/\allowbreak{}api/\allowbreak{}map/\allowbreak{}\{slug\}/\allowbreak{}editability} & which columns a player may edit and what makes each one editable, plus what the pass has to say about the result (EZ1). \\
\id{GET} & \id{/\allowbreak{}api/\allowbreak{}map/\allowbreak{}\{slug\}/\allowbreak{}export} & the Configure export action. For a sketch-originated map (one with a stored sketch layout) it returns a ZIP holding map.xml, level.dat and region/.mca at its top --- a world directory's own contents, which is what a server is handed --- a\,\textcolor{ink30}{\ldots} \\
\id{GET} & \id{/\allowbreak{}api/\allowbreak{}map/\allowbreak{}\{slug\}/\allowbreak{}findings} & everything wrong with this map right now, as far as its stored documents can say, plus the gates a read cannot reach and where to ask them. The operation is MapFindings; this is the door to it. \\
\id{GET} & \id{/\allowbreak{}api/\allowbreak{}map/\allowbreak{}\{slug\}/\allowbreak{}incline} & How steeply the ground is inclined, cell by cell, as characters --- the angle a slope band is picked by (TP24). A cell's glyph is the TENS of its angle, so \id{0} is ground under ten degrees from level, \id{4} is forty to fifty and \id{8} is a face\,\textcolor{ink30}{\ldots} \\
\id{GET} & \id{/\allowbreak{}api/\allowbreak{}map/\allowbreak{}\{slug\}/\allowbreak{}intent} & the map's declarative authoring intent (docs/pgm/new-map-authoring.md), empty if none yet. \\
\id{PUT} & \id{/\allowbreak{}api/\allowbreak{}map/\allowbreak{}\{slug\}/\allowbreak{}intent} & store the intent the author edited and regenerate the map from it. Replaces the stored intent wholesale, which is what makes a deletion in Configure stick. \\
\id{PUT} & \id{/\allowbreak{}api/\allowbreak{}map/\allowbreak{}\{slug\}/\allowbreak{}intent/\allowbreak{}from-plan} & store a freshly compiled intent, carrying the map's authored slices onto it first (IntentCarry), then project it exactly as a normal PUT does. The plan owns the map's structure, so a rebuild is meant to replace its teams, spawns, wools and build zones.\,\textcolor{ink30}{\ldots} \\
\id{PATCH} & \id{/\allowbreak{}api/\allowbreak{}map/\allowbreak{}\{slug\}/\allowbreak{}intent/\allowbreak{}rooms/\allowbreak{}\{reference\}} & put one room piece where the sketch just dragged it. The sketch draws a spawn, a wool room and the building inside one as locked annotations projected out of the intent, so a drag on that canvas has to arrive back here or the picture and\,\textcolor{ink30}{\ldots} \\
\id{GET} & \id{/\allowbreak{}api/\allowbreak{}map/\allowbreak{}\{slug\}/\allowbreak{}islands} & the detected island polygons (from the islands\_json artifact). \\
\id{GET} & \id{/\allowbreak{}api/\allowbreak{}map/\allowbreak{}\{slug\}/\allowbreak{}kit-reach} & can a fresh spawn bridge to each wool with only the placeable blocks its spawn kit grants? (budget-aware traversability). \\
\id{PATCH} & \id{/\allowbreak{}api/\allowbreak{}map/\allowbreak{}\{slug\}/\allowbreak{}metadata} & what a map is called, what it states and who wrote it. The operation is MapMetadata; this is the door to it. \\
\id{GET} & \id{/\allowbreak{}api/\allowbreak{}map/\allowbreak{}\{slug\}/\allowbreak{}monument-seat} & whether each wool monument's block can hold the wool won on it: clear of anything standing in it, and with a block on one of its six faces to be placed against. Both faults stop the wool going in, so both are errors. \\
\id{GET} & \id{/\allowbreak{}api/\allowbreak{}map/\allowbreak{}\{slug\}/\allowbreak{}monument-suggestions} & ?box=x0,y0,z0,x1,y1,z1[\&style=pedestal,label,cap] --- score the pre-gathered monument\_candidate rows inside the author's box for the declared style (F9). No world access: loads the candidates and runs MonumentSuggester.Score. box is required (the author marks the monument area); style defaults to Any,Any,Any.\,\textcolor{ink30}{\ldots} \\
\id{PATCH} & \id{/\allowbreak{}api/\allowbreak{}map/\allowbreak{}\{slug\}/\allowbreak{}observer-spawn} & set/replace the observer spawn. \\
\id{DELETE} & \id{/\allowbreak{}api/\allowbreak{}map/\allowbreak{}\{slug\}/\allowbreak{}observer-spawn} & remove the observer spawn. \\
\id{GET} & \id{/\allowbreak{}api/\allowbreak{}map/\allowbreak{}\{slug\}/\allowbreak{}origin} & \{ sketch: bool \}: whether the map originated in the sketch tool (has a stored sketch layout). The Configure wizard reads this to auto-wire the monument step away for sketch-origin maps (their monuments are derived at export, not authored). \\
\id{GET} & \id{/\allowbreak{}api/\allowbreak{}map/\allowbreak{}\{slug\}/\allowbreak{}plan} & the stored plan blob for a plan-stage map, or \{\} if none. \\
\id{PUT} & \id{/\allowbreak{}api/\allowbreak{}map/\allowbreak{}\{slug\}/\allowbreak{}plan} & replace the stored plan blob (the plan editor's saved state). \\
\id{GET} & \id{/\allowbreak{}api/\allowbreak{}map/\allowbreak{}\{slug\}/\allowbreak{}plan/\allowbreak{}ascii} & the stored plan as a grid of characters, one per proxy cell. The map-scoped twin of GET /plans/\{id\}/ascii: the same render, reached by the slug an authoring driver already holds rather than by a candidate id. ?every=N downsamples. 404 when the map or its plan is missing, 422 when the stored document cannot be read.\,\textcolor{ink30}{\ldots} \\
\id{GET} & \id{/\allowbreak{}api/\allowbreak{}map/\allowbreak{}\{slug\}/\allowbreak{}plan/\allowbreak{}flow} & what the board asks of the two sides, and what that leaves unused, in prose. The companion to the coverage picture and meant to be read before it: a render says where ground is dead and only the flow says why.\,\textcolor{ink30}{\ldots} \\
\id{GET} & \id{/\allowbreak{}api/\allowbreak{}map/\allowbreak{}\{slug\}/\allowbreak{}preflight} & the Review phase's pre-flight gate (new-map-authoring.md \S{}6). Runs the four generated-map checks and reports the export verdict: Round-trip + Mirror --- pure codec/categorizer checks (Preflight).\,\textcolor{ink30}{\ldots} \\
\id{GET} & \id{/\allowbreak{}api/\allowbreak{}map/\allowbreak{}\{slug\}/\allowbreak{}reach} & Which standing ground no player can get to, and why --- the patches, their sizes, their lowest course and the box to stand in. It is the traversability picture's own partition read as numbers: the navigable components, minus the one the board\,\textcolor{ink30}{\ldots} \\
\id{GET} & \id{/\allowbreak{}api/\allowbreak{}map/\allowbreak{}\{slug\}/\allowbreak{}regions} & derived region facets + category counts. \\
\id{POST} & \id{/\allowbreak{}api/\allowbreak{}map/\allowbreak{}\{slug\}/\allowbreak{}regions} & create a primitive region (optionally tagged with the editor draft\_step so it shows in that activity until it's wired --- E10). \\
\id{POST} & \id{/\allowbreak{}api/\allowbreak{}map/\allowbreak{}\{slug\}/\allowbreak{}regions/\allowbreak{}group} & wrap regions in a compound. \\
\id{GET} & \id{/\allowbreak{}api/\allowbreak{}map/\allowbreak{}\{slug\}/\allowbreak{}regions/\allowbreak{}tree} & category-grouped nested region tree (canvas render input). \\
\id{POST} & \id{/\allowbreak{}api/\allowbreak{}map/\allowbreak{}\{slug\}/\allowbreak{}regions/\allowbreak{}ungroup} & dissolve a compound. \\
\id{PATCH} & \id{/\allowbreak{}api/\allowbreak{}map/\allowbreak{}\{slug\}/\allowbreak{}regions/\allowbreak{}\{regionId\}} & rename / bounds / coords. \\
\id{DELETE} & \id{/\allowbreak{}api/\allowbreak{}map/\allowbreak{}\{slug\}/\allowbreak{}regions/\allowbreak{}\{regionId\}} & delete a region (returns an undo snapshot). \\
\id{POST} & \id{/\allowbreak{}api/\allowbreak{}map/\allowbreak{}\{slug\}/\allowbreak{}regions/\allowbreak{}\{regionId\}/\allowbreak{}counterpart} & create the symmetry counterpart(s) of a region (F3). Body: \{mode, center:\{cx,cz\}?\}; centre falls back to the confirmed symmetry artifact. \\
\id{POST} & \id{/\allowbreak{}api/\allowbreak{}map/\allowbreak{}\{slug\}/\allowbreak{}regions/\allowbreak{}\{regionId\}/\allowbreak{}orbit} & fill the symmetry orbit of a region (F3), i.e. create every counterpart implied by the map's confirmed symmetry (rot\_90 $\rightarrow$ 3 turns, mirror/ rot\_180 $\rightarrow$ 1). Used right after a draw so an authored region appears in all symmetric positions.\,\textcolor{ink30}{\ldots} \\
\id{GET} & \id{/\allowbreak{}api/\allowbreak{}map/\allowbreak{}\{slug\}/\allowbreak{}render/\allowbreak{}heightmap} & Elevation as tone, with contour lines every \id{contour} blocks. The read for whether a relief solved into the shape it was drawn as, and the one that shows a flat pad butted against a hill as the ruled edge it is.\,\textcolor{ink30}{\ldots} \\
\id{GET} & \id{/\allowbreak{}api/\allowbreak{}map/\allowbreak{}\{slug\}/\allowbreak{}render/\allowbreak{}mirror} & The board against its own symmetry: the columns that agree with their image, and the ones that do not. The read for whether a board somebody believes is symmetric actually is. \\
\id{GET} & \id{/\allowbreak{}api/\allowbreak{}map/\allowbreak{}\{slug\}/\allowbreak{}render/\allowbreak{}section} & A vertical cut with a Y scale: \id{axis} which way it runs, \id{from}/\id{to} its extent, \id{at} the other coordinate, \id{depth} how far behind it to project. One of only two reads that keep Y, and the one that shows a layer stack, a roof idiom and a riser's courses.\,\textcolor{ink30}{\ldots} \\
\id{GET} & \id{/\allowbreak{}api/\allowbreak{}map/\allowbreak{}\{slug\}/\allowbreak{}render/\allowbreak{}structures} & The building census by block material, for a world the studio did NOT build --- a scanned map, or a corpus one. MISLEADS: It finds roofs by material and cannot see a town this studio built (\id{B149}): a cottage roofed in stained clay reads as ground.\,\textcolor{ink30}{\ldots} \\
\id{GET} & \id{/\allowbreak{}api/\allowbreak{}map/\allowbreak{}\{slug\}/\allowbreak{}render/\allowbreak{}surface} & The paint, read as the tone families \id{TerrainPalette.\allowbreak{}Families} names --- so a board can be checked against the palette it was authored from, and a whole family taken where two members were meant reads as the noise it is.\,\textcolor{ink30}{\ldots} \\
\id{GET} & \id{/\allowbreak{}api/\allowbreak{}map/\allowbreak{}\{slug\}/\allowbreak{}render/\allowbreak{}topdown} & The board from above, one question per image: \id{subject} draws the terrain alone (\id{ground}), what the build recorded itself placing (\id{structure}), the planting (\id{foliage}), the goals and spawns (\id{objectives}), or all of it (\id{combined})\,\textcolor{ink30}{\ldots} \\
\id{GET} & \id{/\allowbreak{}api/\allowbreak{}map/\allowbreak{}\{slug\}/\allowbreak{}render/\allowbreak{}traversability} & The navigable components, with the spawns and goals drawn on them. The read for whether a board hangs together as one place a player can walk. MISLEADS: Headroom is what a player's body passes through, not air: a flower, a torch, a carpet\,\textcolor{ink30}{\ldots} \\
\id{GET} & \id{/\allowbreak{}api/\allowbreak{}map/\allowbreak{}\{slug\}/\allowbreak{}render/\allowbreak{}walk} & What crossing this board charges, drawn: every passable cell shaded by what reaching it from \id{from} costs, with the route to \id{to} over the top. \id{field} picks which of the walk's answers is shaded --- \id{blocks} a player must place, \id{distance}\,\textcolor{ink30}{\ldots} \\
\id{POST} & \id{/\allowbreak{}api/\allowbreak{}map/\allowbreak{}\{slug\}/\allowbreak{}resources} & iron/gold/diamond blocks (optionally in a drawn rect, body \{ bounds?: \{ min\_x, min\_z, max\_x, max\_z \} \}) + how many a \textless{}renewable\textgreater{} already covers. \\
\id{GET} & \id{/\allowbreak{}api/\allowbreak{}map/\allowbreak{}\{slug\}/\allowbreak{}scan-summary} & per-feature breakdowns for the import brief: wool blocks grouped by colour (with a swatch hex) and resource blocks grouped by type, each ordered by count. \\
\id{POST} & \id{/\allowbreak{}api/\allowbreak{}map/\allowbreak{}\{slug\}/\allowbreak{}scan-world} & scan the map's Minecraft world (\textless{}root\textgreater{}/\textless{}slug\textgreater{}/region) and (re)write its relational feature rows (wools/resources/chests/spawners/segments). The world-import half of the pipeline; the xml half is the importer / map editors. \\
\id{GET} & \id{/\allowbreak{}api/\allowbreak{}map/\allowbreak{}\{slug\}/\allowbreak{}segments} & ?axis=x\textbar{}z --- side-view depth profile (B5). Projects the map's vertical solid segments onto a 2D (primary $\times$ y) grid via SideView; feeds the Build-Regions side-view canvas (C7). Mirrors the reference get\_segments. \\
\id{GET} & \id{/\allowbreak{}api/\allowbreak{}map/\allowbreak{}\{slug\}/\allowbreak{}sketch} & the stored sketch layout (the JS-origin blob), or \{\} if none. \\
\id{PUT} & \id{/\allowbreak{}api/\allowbreak{}map/\allowbreak{}\{slug\}/\allowbreak{}sketch} & replace the stored layout blob (the bridge's getState()). A drawing in progress is stored whatever it says. The board's own geometry is checked here and every finding rides back as a complaint, refusals included: a sketch is the author's\,\textcolor{ink30}{\ldots} \\
\id{GET} & \id{/\allowbreak{}api/\allowbreak{}map/\allowbreak{}\{slug\}/\allowbreak{}sketch/\allowbreak{}biome} & the field every column answers to. 404 where the board states none, which is plains everywhere and what every board that never opened the question exports as. \\
\id{PUT} & \id{/\allowbreak{}api/\allowbreak{}map/\allowbreak{}\{slug\}/\allowbreak{}sketch/\allowbreak{}biome} & which biome each column of the exported world carries. Map-wide and answered per chunk, because a biome's tint is blended across a radius and a region drawn to a finer edge never reaches its own colour there. \\
\id{DELETE} & \id{/\allowbreak{}api/\allowbreak{}map/\allowbreak{}\{slug\}/\allowbreak{}sketch/\allowbreak{}biome} & take the field off the board, which is plains everywhere. \\
\id{POST} & \id{/\allowbreak{}api/\allowbreak{}map/\allowbreak{}\{slug\}/\allowbreak{}sketch/\allowbreak{}columns} & the whole built world as per-column runs, which is what the 3-D preview draws (docs/tools/sketch.md). The body is the live layout, the same as the paint and contour previews take, and the stored intent rides along so the structures a map has stated are standing in the picture.\,\textcolor{ink30}{\ldots} \\
\id{DELETE} & \id{/\allowbreak{}api/\allowbreak{}map/\allowbreak{}\{slug\}/\allowbreak{}sketch/\allowbreak{}discard-if-empty} & drop a still-pristine sketch draft (the row ``New sketch'' creates up front, then abandoned). The client calls this best-effort when it leaves the Sketch tool. Discards only a draft that is genuinely untouched: sketch stage, still carrying\,\textcolor{ink30}{\ldots} \\
\id{POST} & \id{/\allowbreak{}api/\allowbreak{}map/\allowbreak{}\{slug\}/\allowbreak{}sketch/\allowbreak{}dressing} & what the dressing pass would place, run against the posted layout and stopped before anything is written. The body is the live layout, the same as the paint preview and the columns read take.\,\textcolor{ink30}{\ldots} \\
\id{POST} & \id{/\allowbreak{}api/\allowbreak{}map/\allowbreak{}\{slug\}/\allowbreak{}sketch/\allowbreak{}finish} & rasterize the stored layout into the world geometry artifacts (layer.parquet / groups.json / segments) so the draft flows into the Configure wizard. 422 only if the layout rasterizes to no ground at all.\,\textcolor{ink30}{\ldots} \\
\id{PUT} & \id{/\allowbreak{}api/\allowbreak{}map/\allowbreak{}\{slug\}/\allowbreak{}sketch/\allowbreak{}from-plan} & replace the stored layout with one a plan compiled, carrying the map's existing finish onto it (CarryFinish). The plan owns the board; the sketch owns its themes, room shells and dressing, and a plan cannot express any of those --- so the compile path merges where SketchPutEndpoint replaces.\,\textcolor{ink30}{\ldots} \\
\id{GET} & \id{/\allowbreak{}api/\allowbreak{}map/\allowbreak{}\{slug\}/\allowbreak{}sketch/\allowbreak{}groups} & every group the board carries, across all its layers, with whether a relief is keyed over it. The list a caller reads before naming a group to draw into or to state terrain for. \\
\id{GET} & \id{/\allowbreak{}api/\allowbreak{}map/\allowbreak{}\{slug\}/\allowbreak{}sketch/\allowbreak{}layers} & the stack in draw order, each layer with the shapes and groups drawn on it. A layer that named itself keeps its id; one that did not is answered under its position, which is the id every other route addresses it by. \\
\id{GET} & \id{/\allowbreak{}api/\allowbreak{}map/\allowbreak{}\{slug\}/\allowbreak{}sketch/\allowbreak{}layers/\allowbreak{}\{layerId\}} & one layer of the stack, with its shapes and groups. 404 where the id names none. \\
\id{PUT} & \id{/\allowbreak{}api/\allowbreak{}map/\allowbreak{}\{slug\}/\allowbreak{}sketch/\allowbreak{}layers/\allowbreak{}\{layerId\}} & state one layer of the stack: the height its ground starts at, whether it is terrain or a made thing, and how its floors meet the ground. Creates the layer at the end of the stack where the id names none.\,\textcolor{ink30}{\ldots} \\
\id{DELETE} & \id{/\allowbreak{}api/\allowbreak{}map/\allowbreak{}\{slug\}/\allowbreak{}sketch/\allowbreak{}layers/\allowbreak{}\{layerId\}} & take one layer, everything drawn on it and the relief of every group that lived only there off the board. \\
\id{PUT} & \id{/\allowbreak{}api/\allowbreak{}map/\allowbreak{}\{slug\}/\allowbreak{}sketch/\allowbreak{}layers/\allowbreak{}\{layerId\}/\allowbreak{}groups/\allowbreak{}\{groupId\}} & state one group whole: what it is called, whether it is fanned onto the symmetry orbit, and which shapes share its ground. Creates it where the layer carries none under that id.\,\textcolor{ink30}{\ldots} \\
\id{DELETE} & \id{/\allowbreak{}api/\allowbreak{}map/\allowbreak{}\{slug\}/\allowbreak{}sketch/\allowbreak{}layers/\allowbreak{}\{layerId\}/\allowbreak{}groups/\allowbreak{}\{groupId\}} & ungroup. The shapes stay on the layer and are drawn where they were drawn; what goes with the group is the orbit fan and the relief keyed under its id. \\
\id{GET} & \id{/\allowbreak{}api/\allowbreak{}map/\allowbreak{}\{slug\}/\allowbreak{}sketch/\allowbreak{}layers/\allowbreak{}\{layerId\}/\allowbreak{}shapes} & the shapes drawn on one layer, in draw order. \\
\id{POST} & \id{/\allowbreak{}api/\allowbreak{}map/\allowbreak{}\{slug\}/\allowbreak{}sketch/\allowbreak{}layers/\allowbreak{}\{layerId\}/\allowbreak{}shapes} & draw one shape on a layer, answering the id it was given. ?group= names the ground it joins, and a name the layer does not carry yet opens a group. A layer that already groups its shapes takes no shape that names none: the orbit fan and the\,\textcolor{ink30}{\ldots} \\
\id{PUT} & \id{/\allowbreak{}api/\allowbreak{}map/\allowbreak{}\{slug\}/\allowbreak{}sketch/\allowbreak{}map-theme} & which registered theme covers every cell no shape's own scope claims. The body is \{``theme'': ``meadow''\}; a null or absent theme clears it, which paints unthemed stone.\,\textcolor{ink30}{\ldots} \\
\id{POST} & \id{/\allowbreak{}api/\allowbreak{}map/\allowbreak{}\{slug\}/\allowbreak{}sketch/\allowbreak{}paint} & the sketch's terrain paint as a palette-indexed block-pixel payload (xs/zs/palette/color\_idx + bounds), which the client expands into the colors array the block-overlay bitmap path already decodes (docs/world-export/terrain-painting.md TP10).\,\textcolor{ink30}{\ldots} \\
\id{POST} & \id{/\allowbreak{}api/\allowbreak{}map/\allowbreak{}\{slug\}/\allowbreak{}sketch/\allowbreak{}probe-footprint} & whether a ring stands on ground, asked of the rasterised footprint the build actually produces rather than of a model of the coast rebuilt outside the studio. The two disagree by a cell or two, and a cell or two is the whole failure: a lift\,\textcolor{ink30}{\ldots} \\
\id{GET} & \id{/\allowbreak{}api/\allowbreak{}map/\allowbreak{}\{slug\}/\allowbreak{}sketch/\allowbreak{}props} & every placement the map carries, with the recipes they name. The recipes ride with the placements because a placement referencing a key nobody can resolve is not readable on its own: a tree states styleKey and the registry states what that key is made of.\,\textcolor{ink30}{\ldots} \\
\id{POST} & \id{/\allowbreak{}api/\allowbreak{}map/\allowbreak{}\{slug\}/\allowbreak{}sketch/\allowbreak{}props} & place one prop, answering the id it was given. A body stating a free id keeps it; one stating none, or one already taken, is minted \{kind\}-\{n\}. The placement goes on the end, because the pass runs in the order props were placed and an addition has not been placed before anything.\,\textcolor{ink30}{\ldots} \\
\id{PATCH} & \id{/\allowbreak{}api/\allowbreak{}map/\allowbreak{}\{slug\}/\allowbreak{}sketch/\allowbreak{}props/\allowbreak{}\{propId\}} & replace one placement, keeping its position in the pass's order. 404 where the id names no placement. \\
\id{DELETE} & \id{/\allowbreak{}api/\allowbreak{}map/\allowbreak{}\{slug\}/\allowbreak{}sketch/\allowbreak{}props/\allowbreak{}\{propId\}} & take one placement off the board. The recipe it named stays in the registry, since a key is shared by every placement wearing it. \\
\id{GET} & \id{/\allowbreak{}api/\allowbreak{}map/\allowbreak{}\{slug\}/\allowbreak{}sketch/\allowbreak{}relief} & every group's relief, by the group id it is solved over. The key is a group rather than a shape because a relief solved per shape leaves a seam wherever two of them meet and disagree about the height they share. \\
\id{POST} & \id{/\allowbreak{}api/\allowbreak{}map/\allowbreak{}\{slug\}/\allowbreak{}sketch/\allowbreak{}relief} & the contour overlay for whatever relief the posted layout carries, one entry per relief-bearing group: its traced lines, its height range, and its bounds. The body is the live layout, the same as the paint preview takes, so the overlay tracks unsaved edits.\,\textcolor{ink30}{\ldots} \\
\id{POST} & \id{/\allowbreak{}api/\allowbreak{}map/\allowbreak{}\{slug\}/\allowbreak{}sketch/\allowbreak{}relief/\allowbreak{}read} & what the relief a posted layout carries charges, per group. Not a walkability score: a relief that is walkable everywhere is a field rather than a map, and a single number ranks every deliberate barrier as a defect.\,\textcolor{ink30}{\ldots} \\
\id{GET} & \id{/\allowbreak{}api/\allowbreak{}map/\allowbreak{}\{slug\}/\allowbreak{}sketch/\allowbreak{}relief/\allowbreak{}\{groupId\}} & one group's relief. 404 where the layout states none for that group, which is every group whose ground is as flat as its shapes drew it. \\
\id{PUT} & \id{/\allowbreak{}api/\allowbreak{}map/\allowbreak{}\{slug\}/\allowbreak{}sketch/\allowbreak{}relief/\allowbreak{}\{groupId\}} & state one group's interior elevation, replacing whatever that group carried. It does not check that the group exists. A relief is authored against a fusion the compiler produces, and whether the id still names one is the question SK1\,\textcolor{ink30}{\ldots} \\
\id{DELETE} & \id{/\allowbreak{}api/\allowbreak{}map/\allowbreak{}\{slug\}/\allowbreak{}sketch/\allowbreak{}relief/\allowbreak{}\{groupId\}} & take one group's relief off the board, which leaves its ground as flat as the shapes drew it. \\
\id{GET} & \id{/\allowbreak{}api/\allowbreak{}map/\allowbreak{}\{slug\}/\allowbreak{}sketch/\allowbreak{}room-styles} & both shells as the stampers will read them. Resolved rather than raw, because the three states a binding has do not survive being handed back as a snapshot: absent falls back to the built-in shell, null asked for open ground, and a stated style is itself.\,\textcolor{ink30}{\ldots} \\
\id{PUT} & \id{/\allowbreak{}api/\allowbreak{}map/\allowbreak{}\{slug\}/\allowbreak{}sketch/\allowbreak{}room-styles/\allowbreak{}\{part\}} & bind the shell one kind of room is stamped in. part is wool or spawn. A body of literal null is a statement, not an omission: it asks for open ground --- a pad rather than a building over it, which is what a spawn on a plateau the plan already shaped often wants to be.\,\textcolor{ink30}{\ldots} \\
\id{DELETE} & \id{/\allowbreak{}api/\allowbreak{}map/\allowbreak{}\{slug\}/\allowbreak{}sketch/\allowbreak{}room-styles/\allowbreak{}\{part\}} & unbind a shell, which puts that kind of room back to its built-in one. Not the same as binding null, which asks for open ground. \\
\id{POST} & \id{/\allowbreak{}api/\allowbreak{}map/\allowbreak{}\{slug\}/\allowbreak{}sketch/\allowbreak{}seats} & where a prop of a stated kind and footprint may stand, which the five declines only ever answer backwards. The dressing pass refuses a placement for the ground it rests on --- a goal's clearance, the map's keep-out, another prop's claim, a\,\textcolor{ink30}{\ldots} \\
\id{GET} & \id{/\allowbreak{}api/\allowbreak{}map/\allowbreak{}\{slug\}/\allowbreak{}sketch/\allowbreak{}shapes/\allowbreak{}\{shapeId\}} & one shape, wherever on the stack it is drawn. 404 where the id names none. \\
\id{PATCH} & \id{/\allowbreak{}api/\allowbreak{}map/\allowbreak{}\{slug\}/\allowbreak{}sketch/\allowbreak{}shapes/\allowbreak{}\{shapeId\}} & change one shape without restating the board. A stated field replaces what the shape carried and a stated null takes the field off, so the one call both writes a height and clears a relief scope. id is the address and is kept whatever the\,\textcolor{ink30}{\ldots} \\
\id{DELETE} & \id{/\allowbreak{}api/\allowbreak{}map/\allowbreak{}\{slug\}/\allowbreak{}sketch/\allowbreak{}shapes/\allowbreak{}\{shapeId\}} & rub one shape out, and take it out of every group that listed it. \\
\id{POST} & \id{/\allowbreak{}api/\allowbreak{}map/\allowbreak{}\{slug\}/\allowbreak{}sketch/\allowbreak{}shapes/\allowbreak{}\{shapeId\}/\allowbreak{}bend} & redraw one outline as a coast. The compiler emits a staircase of the plan's rectangles, which is the board's shape and not its coast. Redrawing the ring by hand states the coast twice --- once in the plan and once in the sketch, free to disagree --- so the bend is taken over whatever the compile produced instead.\,\textcolor{ink30}{\ldots} \\
\id{POST} & \id{/\allowbreak{}api/\allowbreak{}map/\allowbreak{}\{slug\}/\allowbreak{}sketch/\allowbreak{}shapes/\allowbreak{}\{shapeId\}/\allowbreak{}vertices} & add one point to one outline, after the vertex after names. Stating no x/z puts it at the midpoint of that edge, which is a new corner half way along a wall with nothing else moved; the answer carries the index it landed at, which is the address the move route then takes.\,\textcolor{ink30}{\ldots} \\
\id{PATCH} & \id{/\allowbreak{}api/\allowbreak{}map/\allowbreak{}\{slug\}/\allowbreak{}sketch/\allowbreak{}shapes/\allowbreak{}\{shapeId\}/\allowbreak{}vertices/\allowbreak{}\{index\}} & move one point of one outline. Every other vertex stays exactly where it was drawn. That is the whole of the call: a board's shapes abut, and an edit that drags a ring's other points opens ground between two that were flush, which is what a whole-ring transform does and what a hand pulling one corner does not.\,\textcolor{ink30}{\ldots} \\
\id{DELETE} & \id{/\allowbreak{}api/\allowbreak{}map/\allowbreak{}\{slug\}/\allowbreak{}sketch/\allowbreak{}shapes/\allowbreak{}\{shapeId\}/\allowbreak{}vertices/\allowbreak{}\{index\}} & take one point out of one outline, leaving every other where it was drawn. Refused where the outline is down to its last three, since two points draw no ground, and where dropping the point folds the ring. \\
\id{GET} & \id{/\allowbreak{}api/\allowbreak{}map/\allowbreak{}\{slug\}/\allowbreak{}sketch/\allowbreak{}themes} & the registry, and which of it is the map default. A theme that will not parse as one is left out, the same way the painter drops it: the registry is a snapshot store rather than a validated model, and a reader asking what the board paints with wants what the board will actually paint with.\,\textcolor{ink30}{\ldots} \\
\id{GET} & \id{/\allowbreak{}api/\allowbreak{}map/\allowbreak{}\{slug\}/\allowbreak{}sketch/\allowbreak{}themes/\allowbreak{}\{themeId\}} & one theme, as the painter reads it. 404 where the registry carries no such id, or carries one that will not parse as a theme. \\
\id{PUT} & \id{/\allowbreak{}api/\allowbreak{}map/\allowbreak{}\{slug\}/\allowbreak{}sketch/\allowbreak{}themes/\allowbreak{}\{themeId\}} & register one theme under an id, replacing whatever that id carried. A registry entry is addressed by the name an author gave it, so this is the one write in the sketch that creates and replaces through the same verb. \\
\id{DELETE} & \id{/\allowbreak{}api/\allowbreak{}map/\allowbreak{}\{slug\}/\allowbreak{}sketch/\allowbreak{}themes/\allowbreak{}\{themeId\}} & take a theme out of the registry. It does not refuse over what still names the id. A shape painting with a theme the registry has stopped carrying takes the map default, and the map default naming one takes unthemed stone --- both are already\,\textcolor{ink30}{\ldots} \\
\id{GET} & \id{/\allowbreak{}api/\allowbreak{}map/\allowbreak{}\{slug\}/\allowbreak{}slopes} & The worst step to a neighbour per sampled cell, classed \id{.\allowbreak{}} walked, \id{:\allowbreak{}} scrambled with a block, \id{\#} a barrier --- the same tiers a walk is priced in. A cliff reads as a line, a ramp as a band crossing it, and an overdone relief as a page.\,\textcolor{ink30}{\ldots} \\
\id{POST} & \id{/\allowbreak{}api/\allowbreak{}map/\allowbreak{}\{slug\}/\allowbreak{}spawns} & link a spawn to a region. \\
\id{PATCH} & \id{/\allowbreak{}api/\allowbreak{}map/\allowbreak{}\{slug\}/\allowbreak{}spawns/\allowbreak{}\{regionId\}} & update a spawn link. \\
\id{DELETE} & \id{/\allowbreak{}api/\allowbreak{}map/\allowbreak{}\{slug\}/\allowbreak{}spawns/\allowbreak{}\{regionId\}} & remove a spawn link. \\
\id{GET} & \id{/\allowbreak{}api/\allowbreak{}map/\allowbreak{}\{slug\}/\allowbreak{}state} & where this map has got to, which authoring artifacts it holds, and what may be done to it from here. A tool asks about the map it has open so it can tell an origination from a rebuild before offering the action rather than after performing it; the moves are the same question answered outright.\,\textcolor{ink30}{\ldots} \\
\id{GET} & \id{/\allowbreak{}api/\allowbreak{}map/\allowbreak{}\{slug\}/\allowbreak{}stroke} & A drawn stroke walked end to end, down its own centreline rather than between its points: per block the ground, whether the paving reaches it, what it is made of and the step from the block before, then the worst step, the material runs and every stretch the paving does not reach.\,\textcolor{ink30}{\ldots} \\
\id{GET} & \id{/\allowbreak{}api/\allowbreak{}map/\allowbreak{}\{slug\}/\allowbreak{}symmetry} & global symmetry of the map's islands (B7). Returns the cached symmetry\_json artifact, or computes it on demand from the islands\_json artifact (excluding the Configure-excluded islands) and caches it with status ``unconfirmed''. \\
\id{PATCH} & \id{/\allowbreak{}api/\allowbreak{}map/\allowbreak{}\{slug\}/\allowbreak{}symmetry} & confirm/reject the detected symmetry (B7). Updates status (``confirmed''/``none''), an optional user-override confirmed\_type, and an optional centre override. Mirrors the reference patch\_symmetry. \\
\id{POST} & \id{/\allowbreak{}api/\allowbreak{}map/\allowbreak{}\{slug\}/\allowbreak{}teams} & add a team. \\
\id{PATCH} & \id{/\allowbreak{}api/\allowbreak{}map/\allowbreak{}\{slug\}/\allowbreak{}teams/\allowbreak{}\{teamId\}} & update a team. \\
\id{DELETE} & \id{/\allowbreak{}api/\allowbreak{}map/\allowbreak{}\{slug\}/\allowbreak{}teams/\allowbreak{}\{teamId\}} & remove a team and its spawns. \\
\id{GET} & \id{/\allowbreak{}api/\allowbreak{}map/\allowbreak{}\{slug\}/\allowbreak{}themes/\allowbreak{}census} & Every ground cell counted by the theme that paints it: cells per theme and their share, the distinct surface materials each carries, which theme borders which and over how many cells, and the board's whole palette count.\,\textcolor{ink30}{\ldots} \\
\id{GET} & \id{/\allowbreak{}api/\allowbreak{}map/\allowbreak{}\{slug\}/\allowbreak{}top-surface} & per-column surface colour overlay (B4). Reads the cached layer.parquet artifact, maps each column's (block\_id, block\_data) to a hex colour, and returns parallel xs/zs/colors arrays + the bounds.\,\textcolor{ink30}{\ldots} \\
\id{GET} & \id{/\allowbreak{}api/\allowbreak{}map/\allowbreak{}\{slug\}/\allowbreak{}transect} & A polyline walked block by block, answered as numbers: \id{points} is the line as \id{x,z;x,z[;x,z.\allowbreak{}.\allowbreak{}.\allowbreak{}]}, \id{every} thins the stations, \id{beside} lists every distinct claim within that many cells of the line.\,\textcolor{ink30}{\ldots} \\
\id{GET} & \id{/\allowbreak{}api/\allowbreak{}map/\allowbreak{}\{slug\}/\allowbreak{}traversability} & spawn$\leftrightarrow$wool connectivity. \\
\id{GET} & \id{/\allowbreak{}api/\allowbreak{}map/\allowbreak{}\{slug\}/\allowbreak{}walk} & The same journey as numbers rather than as a picture: whether the place can be reached, how far it is in blocks, how many blocks a player must place to get there --- the climb at a rise of delta costing delta minus one, and one a cell for void bridged --- and how many falls over the free height it takes.\,\textcolor{ink30}{\ldots} \\
\id{GET} & \id{/\allowbreak{}api/\allowbreak{}map/\allowbreak{}\{slug\}/\allowbreak{}wool-availability} & per declared wool, is it obtainable? \\
\id{POST} & \id{/\allowbreak{}api/\allowbreak{}map/\allowbreak{}\{slug\}/\allowbreak{}wool-sources} & wool colours found inside a drawn rectangle (body: \{ bounds: \{ min\_x, min\_z, max\_x, max\_z \} \}). \\
\id{GET} & \id{/\allowbreak{}api/\allowbreak{}map/\allowbreak{}\{slug\}/\allowbreak{}wool-suggestions} & wool colours in the world not yet declared as objectives. \\
\id{POST} & \id{/\allowbreak{}api/\allowbreak{}map/\allowbreak{}\{slug\}/\allowbreak{}wools} & add a wool objective. \\
\id{PATCH} & \id{/\allowbreak{}api/\allowbreak{}map/\allowbreak{}\{slug\}/\allowbreak{}wools/\allowbreak{}\{woolId\}} & update a wool. \\
\id{DELETE} & \id{/\allowbreak{}api/\allowbreak{}map/\allowbreak{}\{slug\}/\allowbreak{}wools/\allowbreak{}\{woolId\}} & remove a wool. \\
\id{POST} & \id{/\allowbreak{}api/\allowbreak{}map/\allowbreak{}\{slug\}/\allowbreak{}wools/\allowbreak{}\{woolId\}/\allowbreak{}monuments} & add a monument to a wool. \\
\id{PATCH} & \id{/\allowbreak{}api/\allowbreak{}map/\allowbreak{}\{slug\}/\allowbreak{}wools/\allowbreak{}\{woolId\}/\allowbreak{}monuments/\allowbreak{}\{monId\}} & update a monument. \\
\id{DELETE} & \id{/\allowbreak{}api/\allowbreak{}map/\allowbreak{}\{slug\}/\allowbreak{}wools/\allowbreak{}\{woolId\}/\allowbreak{}monuments/\allowbreak{}\{monId\}} & remove a monument. \\
\id{GET} & \id{/\allowbreak{}api/\allowbreak{}map/\allowbreak{}\{slug\}/\allowbreak{}xml} & render the map as a PGM map.xml (proto 1.5.0). The export leg of the round-trip: doc $\rightarrow$ Deserializer.FromDict $\rightarrow$ MapXml $\rightarrow$ XmlWriter.ToXml --- the same path the round-trip harness (check \#2) exercises across the corpus.\,\textcolor{ink30}{\ldots} \\
\addlinespace[3pt]
\multicolumn{3}{@{}l}{\textbf{\id{/terrain}} \textcolor{ink60}{\small(12)}}\\[1pt]
\midrule[0.2pt]
\id{GET} & \id{/\allowbreak{}api/\allowbreak{}terrain/\allowbreak{}biomes} & the biomes a field may name, each with the grass colour it tints ground with (docs/world-export/terrain-painting.md \S{}5b). Served rather than restated in the client for the reason the block palette is: the ids live in Biome and the colours\,\textcolor{ink30}{\ldots} \\
\id{GET} & \id{/\allowbreak{}api/\allowbreak{}terrain/\allowbreak{}blocks} & the blocks a terrain-paint material may be built from (TerrainPalette), each with its id/data pair, display name, picker group and swatch colour. The block picker reads it, so the colour a picker shows is the colour the export places (docs/world-export/terrain-painting.md \S{}3).\,\textcolor{ink30}{\ldots} \\
\id{GET} & \id{/\allowbreak{}api/\allowbreak{}terrain/\allowbreak{}boulder-forms} & the four rock shapes, each an actual rock. \\
\id{GET} & \id{/\allowbreak{}api/\allowbreak{}terrain/\allowbreak{}looks} & the construction words a block's construction is drawn from, each with what it means. Served rather than restated by every caller, for the same reason the block palette is: a second copy on the far side of the wire is a vocabulary free to drift from the one the table was measured under.\,\textcolor{ink30}{\ldots} \\
\id{POST} & \id{/\allowbreak{}api/\allowbreak{}terrain/\allowbreak{}material-preview} & body is one serialized TerrainMaterial; returns both views of it (top-down and cut open). What a style editor re-renders on every edit --- see StylePreview for why one material needs two views. ?format=png\&view=plan\textbar{}section answers that one view as PNG bytes instead.\,\textcolor{ink30}{\ldots} \\
\id{GET} & \id{/\allowbreak{}api/\allowbreak{}terrain/\allowbreak{}patterns} & every material kind a theme may be built from, with what it does, the facts about a cell it varies with, and its fields with their types and defaults (MaterialVocabulary).\,\textcolor{ink30}{\ldots} \\
\id{POST} & \id{/\allowbreak{}api/\allowbreak{}terrain/\allowbreak{}prop-preview} & one placed prop and the terrain finish it stands on; returns a sample patch the pass actually dressed, from above and cut open. The theme is part of the request because what the paint leaves on top is what decides whether flora grows at all and what a path may repaint.\,\textcolor{ink30}{\ldots} \\
\id{GET} & \id{/\allowbreak{}api/\allowbreak{}terrain/\allowbreak{}species} & every vanilla species, each built. \\
\id{GET} & \id{/\allowbreak{}api/\allowbreak{}terrain/\allowbreak{}stroke-styles} & the five ways a stroke paves the ground it crosses, each drawn. \\
\id{POST} & \id{/\allowbreak{}api/\allowbreak{}terrain/\allowbreak{}theme-map-preview} & body is a plan JSON; compiles it, paints the terrain through the scoped theme resolver, and returns a top-down SVG of the map's top blocks so the Theme rail's apply step can show the themes on the actual map (TP10). 400 when the plan can't be rendered.\,\textcolor{ink30}{\ldots} \\
\id{POST} & \id{/\allowbreak{}api/\allowbreak{}terrain/\allowbreak{}theme-preview} & body is a serialized terrain-paint theme (a TerrainTheme JSON); returns the sample plateau painted with it and cut open, plus a top-down swatch per themeable bucket, so a theme editor previews the whole finish and each brush as it is edited (TP10).\,\textcolor{ink30}{\ldots} \\
\id{GET} & \id{/\allowbreak{}api/\allowbreak{}terrain/\allowbreak{}water-forms} & the three channel forms, each an actual dug channel seen from above. \\
\addlinespace[3pt]
\multicolumn{3}{@{}l}{\textbf{\id{/room-styles}} \textcolor{ink60}{\small(10)}}\\[1pt]
\midrule[0.2pt]
\id{GET} & \id{/\allowbreak{}api/\allowbreak{}room-styles} & the room-style library, newest first, each with the shell it stamps. \\
\id{POST} & \id{/\allowbreak{}api/\allowbreak{}room-styles} & compose a room style from existing styles. 400 \id{\{error, findings\}} when the composed shell fails Check --- a block named for a geometric role that is not that kind of block, a doorway that does not clear the least height a door may, or a roof whose own materials are wrong for its pitch or its family.\,\textcolor{ink30}{\ldots} \\
\id{GET} & \id{/\allowbreak{}api/\allowbreak{}room-styles/\allowbreak{}block-kinds} & what kind of block each house-style field takes, and the ids of each kind. Served for the reason the door list is: the authoritative statement is HouseBlockKinds in PgmStudio.Minecraft, which neither the client nor an agent can reach, and\,\textcolor{ink30}{\ldots} \\
\id{GET} & \id{/\allowbreak{}api/\allowbreak{}room-styles/\allowbreak{}doors} & the doors a room may be stamped with. Served rather than restated in the client: the authoritative list is the same DoorMaterials table the wool-room block filter is built from, so a door offered here is always one an attacker can break. \\
\id{POST} & \id{/\allowbreak{}api/\allowbreak{}room-styles/\allowbreak{}preview} & the shell a set of courses composes to, previewed without saving any of it. What the editor re-renders as courses are bound and knobs are turned. \\
\id{POST} & \id{/\allowbreak{}api/\allowbreak{}room-styles/\allowbreak{}preview-snapshot} & body is a serialized HouseStyle; returns the shell it stamps. What a map's bound style is previewed through: the binding is a snapshot rather than a library id (structures.md \S{}9), so the picture has to come from the snapshot too --- reading\,\textcolor{ink30}{\ldots} \\
\id{GET} & \id{/\allowbreak{}api/\allowbreak{}room-styles/\allowbreak{}\{id\}} & one room style with its per-part courses. \\
\id{PUT} & \id{/\allowbreak{}api/\allowbreak{}room-styles/\allowbreak{}\{id\}} & replace a room style's knobs and its whole set of courses. Refuses the same way RoomStyleCreateEndpoint does. \\
\id{DELETE} & \id{/\allowbreak{}api/\allowbreak{}room-styles/\allowbreak{}\{id\}} & forget a room style (its courses cascade; the styles stay). \\
\id{GET} & \id{/\allowbreak{}api/\allowbreak{}room-styles/\allowbreak{}\{id\}/\allowbreak{}json} & the room style assembled into the stamper's own JSON: the form the export consumes and the form a map snapshots when it binds one. \\
\addlinespace[3pt]
\multicolumn{3}{@{}l}{\textbf{\id{/plan}} \textcolor{ink60}{\small(9)}}\\[1pt]
\midrule[0.2pt]
\id{POST} & \id{/\allowbreak{}api/\allowbreak{}plan} & originate a blank authored plan: create a map row at stage=plan with an empty plan\_json artifact (no candidate provenance). Returns the slug; the client navigates to /maps/\{slug\}/plan, where the editor keeps its default blank document until first save.\,\textcolor{ink30}{\ldots} \\
\id{POST} & \id{/\allowbreak{}api/\allowbreak{}plan/\allowbreak{}ascii} & the board as a grid of characters, one per proxy cell, off a posted plan. The read that shows a relation between two rectangles. The other projections answer numbers and geometry, and a number cannot say that a sixteen-cell bar is reached\,\textcolor{ink30}{\ldots} \\
\id{POST} & \id{/\allowbreak{}api/\allowbreak{}plan/\allowbreak{}columns} & the world a plan builds, as per-column runs, which is what the plan tool's 3-D preview draws (docs/tools/plan.md). The body is the plan document itself. A compiled plan carries a full intent, so unlike the sketch side this shows the wool\,\textcolor{ink30}{\ldots} \\
\id{POST} & \id{/\allowbreak{}api/\allowbreak{}plan/\allowbreak{}compile} & compile a plan wire document one-way into the pair the draft pipeline consumes: \{ layout: \textless{}SketchLayout\textgreater{}, intent: \textless{}MapIntent\textgreater{} \}. Structural validator errors block the compile (422, carrying the error findings); lint alone never blocks.\,\textcolor{ink30}{\ldots} \\
\id{POST} & \id{/\allowbreak{}api/\allowbreak{}plan/\allowbreak{}evaluate} & the plan editor's live rule-evaluator score + lint feed (the critic that scores a .plan.json). The request body is a plan wire document; the response is an EvaluationDto: the summed score (lower is better, 0 perfect), a valid flag (no hard\,\textcolor{ink30}{\ldots} \\
\id{POST} & \id{/\allowbreak{}api/\allowbreak{}plan/\allowbreak{}feasibility} & a diagnostic about the composer, not a read about the board. Could the composer have produced this plan? The question PlanEvaluateEndpoint does not ask, and the two genuinely diverge --- a plan can score 0 with no rule fired and still be outside everything the emitters can build.\,\textcolor{ink30}{\ldots} \\
\id{POST} & \id{/\allowbreak{}api/\allowbreak{}plan/\allowbreak{}inspect} & the live derived-geometry feed for the plan editor's canvas overlays. The request body is a plan wire document (.plan.json); the response carries everything already resolved to block coordinates so the canvas draws it directly: interfaces\,\textcolor{ink30}{\ldots} \\
\id{POST} & \id{/\allowbreak{}api/\allowbreak{}plan/\allowbreak{}room} & ?piece=\textless{}id\textgreater{} --- the contents a spawn or wool-room piece carries: the marker the room is built around, the footprint the building stands on, and the spawn's iron cube.\,\textcolor{ink30}{\ldots} \\
\id{POST} & \id{/\allowbreak{}api/\allowbreak{}plan/\allowbreak{}\{planId\}/\allowbreak{}author} & commit a generator plan candidate to authoring: create a map row at stage=plan seeded with the candidate's plan\_json (a plan\_json artifact) and a plan\_source\_id back to the candidate.\,\textcolor{ink30}{\ldots} \\
\addlinespace[3pt]
\multicolumn{3}{@{}l}{\textbf{\id{/themes}} \textcolor{ink60}{\small(8)}}\\[1pt]
\midrule[0.2pt]
\id{GET} & \id{/\allowbreak{}api/\allowbreak{}themes} & the theme library, newest first, each with the sample plateau it finishes. \\
\id{POST} & \id{/\allowbreak{}api/\allowbreak{}themes} & compose a theme from existing styles (the knobs + bucket$\rightarrow$style bindings). \\
\id{POST} & \id{/\allowbreak{}api/\allowbreak{}themes/\allowbreak{}import} & lift a whole theme JSON into the library: one style per bucket + a theme binding them. 400, never 500, on invalid theme JSON. \\
\id{POST} & \id{/\allowbreak{}api/\allowbreak{}themes/\allowbreak{}preview} & the theme a set of bindings composes to, previewed without saving any of it. What the library's theme editor re-renders as buckets are bound and knobs are turned. \\
\id{GET} & \id{/\allowbreak{}api/\allowbreak{}themes/\allowbreak{}\{id\}} & a theme with its per-bucket style bindings. \\
\id{PUT} & \id{/\allowbreak{}api/\allowbreak{}themes/\allowbreak{}\{id\}} & replace a theme's knobs and its whole set of bucket bindings. \\
\id{DELETE} & \id{/\allowbreak{}api/\allowbreak{}themes/\allowbreak{}\{id\}} & forget a theme (its bucket bindings cascade; the styles stay). \\
\id{GET} & \id{/\allowbreak{}api/\allowbreak{}themes/\allowbreak{}\{id\}/\allowbreak{}json} & the theme assembled into the painter's theme JSON (the form the export consumes and a map snapshots when it applies the theme). \\
\addlinespace[3pt]
\multicolumn{3}{@{}l}{\textbf{\id{/boulder-styles}} \textcolor{ink60}{\small(7)}}\\[1pt]
\midrule[0.2pt]
\id{GET} & \id{/\allowbreak{}api/\allowbreak{}boulder-styles} & \textcolor{ink30}{---} \\
\id{POST} & \id{/\allowbreak{}api/\allowbreak{}boulder-styles} & \textcolor{ink30}{---} \\
\id{POST} & \id{/\allowbreak{}api/\allowbreak{}boulder-styles/\allowbreak{}preview} & \textcolor{ink30}{---} \\
\id{GET} & \id{/\allowbreak{}api/\allowbreak{}boulder-styles/\allowbreak{}\{id\}} & \textcolor{ink30}{---} \\
\id{PUT} & \id{/\allowbreak{}api/\allowbreak{}boulder-styles/\allowbreak{}\{id\}} & \textcolor{ink30}{---} \\
\id{DELETE} & \id{/\allowbreak{}api/\allowbreak{}boulder-styles/\allowbreak{}\{id\}} & \textcolor{ink30}{---} \\
\id{GET} & \id{/\allowbreak{}api/\allowbreak{}boulder-styles/\allowbreak{}\{id\}/\allowbreak{}json} & \textcolor{ink30}{---} \\
\addlinespace[3pt]
\multicolumn{3}{@{}l}{\textbf{\id{/plans}} \textcolor{ink60}{\small(7)}}\\[1pt]
\midrule[0.2pt]
\id{GET} & \id{/\allowbreak{}api/\allowbreak{}plans} & [?origin=generated\textbar{}authored\textbar{}imported] --- the open-from-DB browser list, newest touched first. Summaries only (no plan JSON); the detail endpoint carries the document. \\
\id{POST} & \id{/\allowbreak{}api/\allowbreak{}plans} & save the plan open in the editor. Applies the fork-or-mutate doctrine in SaveFromEditorAsync and returns the resulting row. A malformed plan body is answered 400, never 500. \\
\id{GET} & \id{/\allowbreak{}api/\allowbreak{}plans/\allowbreak{}\{id\}} & one plan with its *.plan.json document, to load into the editor. \\
\id{DELETE} & \id{/\allowbreak{}api/\allowbreak{}plans/\allowbreak{}\{id\}} & forget a plan (204). Forks of it survive: the self-FK sets their parent\_id null rather than cascading. \\
\id{GET} & \id{/\allowbreak{}api/\allowbreak{}plans/\allowbreak{}\{id\}/\allowbreak{}ascii} & the same fanned board as a grid of characters, one per proxy cell. A plan is a list of rectangles measured in cells, and most of what goes wrong with one is a relation between two of them --- a landform wider than the band that reaches it, a wall on the only throat.\,\textcolor{ink30}{\ldots} \\
\id{GET} & \id{/\allowbreak{}api/\allowbreak{}plans/\allowbreak{}\{id\}/\allowbreak{}png} & the same fanned board as PlanSvgEndpoint, rasterized: a picture an image reader can actually open, where the vector card cannot be. Draws off PlanBoardScene, the geometry the SVG endpoint shares, so the two can never disagree about what the plan looks like. 404 when the plan is missing.\,\textcolor{ink30}{\ldots} \\
\id{GET} & \id{/\allowbreak{}api/\allowbreak{}plans/\allowbreak{}\{id\}/\allowbreak{}svg} & render a stored plan to the same board SVG the browse feed uses, so the hold tray can show a thumbnail of a persisted plan. 404 when the plan is missing. \\
\addlinespace[3pt]
\multicolumn{3}{@{}l}{\textbf{\id{/tree-styles}} \textcolor{ink60}{\small(7)}}\\[1pt]
\midrule[0.2pt]
\id{GET} & \id{/\allowbreak{}api/\allowbreak{}tree-styles} & the tree library, newest first, each drawn through the grower the export runs. A tree is picked by what it looks like: six woods differ in colour and six species in shape, and neither reads off a number. \\
\id{POST} & \id{/\allowbreak{}api/\allowbreak{}tree-styles} & \textcolor{ink30}{---} \\
\id{POST} & \id{/\allowbreak{}api/\allowbreak{}tree-styles/\allowbreak{}preview} & \textcolor{ink30}{---} \\
\id{GET} & \id{/\allowbreak{}api/\allowbreak{}tree-styles/\allowbreak{}\{id\}} & \textcolor{ink30}{---} \\
\id{PUT} & \id{/\allowbreak{}api/\allowbreak{}tree-styles/\allowbreak{}\{id\}} & \textcolor{ink30}{---} \\
\id{DELETE} & \id{/\allowbreak{}api/\allowbreak{}tree-styles/\allowbreak{}\{id\}} & No question is asked, because nothing binds a recipe: a placement names a key in its own document's registry, which the pull copied, so a map keeps its trees when the row they were pulled from goes. \\
\id{GET} & \id{/\allowbreak{}api/\allowbreak{}tree-styles/\allowbreak{}\{id\}/\allowbreak{}json} & the recipe as the dressing document states it, which is what a pull copies into a map's registry under a key. \\
\addlinespace[3pt]
\multicolumn{3}{@{}l}{\textbf{\id{/biome-patterns}} \textcolor{ink60}{\small(6)}}\\[1pt]
\midrule[0.2pt]
\id{GET} & \id{/\allowbreak{}api/\allowbreak{}biome-patterns} & [?kind=solid\textbar{}cell\textbar{}noise] --- the biome library, newest first, each with the patch of ground it tints. \\
\id{POST} & \id{/\allowbreak{}api/\allowbreak{}biome-patterns} & save a new named field. \\
\id{POST} & \id{/\allowbreak{}api/\allowbreak{}biome-patterns/\allowbreak{}preview} & what a draft field draws, saving nothing. Body is a bare BiomeField, unwrapped, the way the material preview takes a bare material. \\
\id{GET} & \id{/\allowbreak{}api/\allowbreak{}biome-patterns/\allowbreak{}\{id\}} & one pattern. \\
\id{PUT} & \id{/\allowbreak{}api/\allowbreak{}biome-patterns/\allowbreak{}\{id\}} & update a pattern in place. A map holds a snapshot rather than a key into this library, so an edit here retints nothing already built. \\
\id{DELETE} & \id{/\allowbreak{}api/\allowbreak{}biome-patterns/\allowbreak{}\{id\}} & forget a pattern. Nothing is asked first: a map took a copy, so no board changes when the row goes. \\
\addlinespace[3pt]
\multicolumn{3}{@{}l}{\textbf{\id{/porch-styles}} \textcolor{ink60}{\small(6)}}\\[1pt]
\midrule[0.2pt]
\id{GET} & \id{/\allowbreak{}api/\allowbreak{}porch-styles} & the porch library, newest first, each fronting the sample building. \\
\id{POST} & \id{/\allowbreak{}api/\allowbreak{}porch-styles} & \textcolor{ink30}{---} \\
\id{POST} & \id{/\allowbreak{}api/\allowbreak{}porch-styles/\allowbreak{}preview} & \textcolor{ink30}{---} \\
\id{GET} & \id{/\allowbreak{}api/\allowbreak{}porch-styles/\allowbreak{}\{id\}} & \textcolor{ink30}{---} \\
\id{PUT} & \id{/\allowbreak{}api/\allowbreak{}porch-styles/\allowbreak{}\{id\}} & \textcolor{ink30}{---} \\
\id{DELETE} & \id{/\allowbreak{}api/\allowbreak{}porch-styles/\allowbreak{}\{id\}} & \textcolor{ink30}{---} \\
\addlinespace[3pt]
\multicolumn{3}{@{}l}{\textbf{\id{/roof-styles}} \textcolor{ink60}{\small(6)}}\\[1pt]
\midrule[0.2pt]
\id{GET} & \id{/\allowbreak{}api/\allowbreak{}roof-styles} & the roof library, newest first, each on the sample building it tops. \\
\id{POST} & \id{/\allowbreak{}api/\allowbreak{}roof-styles} & 400 \id{\{error, findings\}} on anything CheckRoof names --- a roof or a verge that is a log or a ground material, a slab block that is not a slab, and a slab named as a whole-block roof.\,\textcolor{ink30}{\ldots} \\
\id{POST} & \id{/\allowbreak{}api/\allowbreak{}roof-styles/\allowbreak{}preview} & the building a draft roof tops, previewed without saving it. \\
\id{GET} & \id{/\allowbreak{}api/\allowbreak{}roof-styles/\allowbreak{}\{id\}} & one roof with its per-part courses. \\
\id{PUT} & \id{/\allowbreak{}api/\allowbreak{}roof-styles/\allowbreak{}\{id\}} & Refuses the same way RoofStyleCreateEndpoint does. \\
\id{DELETE} & \id{/\allowbreak{}api/\allowbreak{}roof-styles/\allowbreak{}\{id\}} & forget a roof, unless a house still wears it. A part bound by a building is one the library must refuse to forget, the answer a style already gives a theme. \\
\addlinespace[3pt]
\multicolumn{3}{@{}l}{\textbf{\id{/storey-styles}} \textcolor{ink60}{\small(6)}}\\[1pt]
\midrule[0.2pt]
\id{GET} & \id{/\allowbreak{}api/\allowbreak{}storey-styles} & the storey library, newest first, each as the one-storey building it makes. \\
\id{POST} & \id{/\allowbreak{}api/\allowbreak{}storey-styles} & 400 \id{\{error, findings\}} when the storey's window names a block that is not the kind its form needs --- a stair lattice or an arched window whose block is not a stair, or a slab band whose block is not a single slab. \\
\id{POST} & \id{/\allowbreak{}api/\allowbreak{}storey-styles/\allowbreak{}preview} & \textcolor{ink30}{---} \\
\id{GET} & \id{/\allowbreak{}api/\allowbreak{}storey-styles/\allowbreak{}\{id\}} & \textcolor{ink30}{---} \\
\id{PUT} & \id{/\allowbreak{}api/\allowbreak{}storey-styles/\allowbreak{}\{id\}} & Refuses the same way StoreyStyleCreateEndpoint does. \\
\id{DELETE} & \id{/\allowbreak{}api/\allowbreak{}storey-styles/\allowbreak{}\{id\}} & \textcolor{ink30}{---} \\
\addlinespace[3pt]
\multicolumn{3}{@{}l}{\textbf{\id{/styles}} \textcolor{ink60}{\small(5)}}\\[1pt]
\midrule[0.2pt]
\id{GET} & \id{/\allowbreak{}api/\allowbreak{}styles} & [?kind=voronoi\textbar{}noise\textbar{}\ldots{}] --- the style library, newest first, optionally one kind (the ``show every voronoi'' browse). \\
\id{POST} & \id{/\allowbreak{}api/\allowbreak{}styles} & save a new reusable style. \\
\id{GET} & \id{/\allowbreak{}api/\allowbreak{}styles/\allowbreak{}\{id\}} & one style. \\
\id{PUT} & \id{/\allowbreak{}api/\allowbreak{}styles/\allowbreak{}\{id\}} & update a style in place (edits every theme that binds it --- a library edit, not a map's applied snapshot). \\
\id{DELETE} & \id{/\allowbreak{}api/\allowbreak{}styles/\allowbreak{}\{id\}} & forget a style. Refused with 409 and the names of the themes and room styles still binding it, since a style is shared by both libraries and its bindings are what the foreign key would otherwise complain about. \\
\addlinespace[3pt]
\multicolumn{3}{@{}l}{\textbf{\id{/maps}} \textcolor{ink60}{\small(3)}}\\[1pt]
\midrule[0.2pt]
\id{GET} & \id{/\allowbreak{}api/\allowbreak{}maps} & [?stage=plan\textbar{}sketch\textbar{}configure\textbar{}edit] --- the dashboard map list. stage selects a collection, and the four collections are not all the same kind of thing, because a map's layers and its progress are not the same thing. plan and sketch name\,\textcolor{ink30}{\ldots} \\
\id{GET} & \id{/\allowbreak{}api/\allowbreak{}maps/\allowbreak{}import-candidates} & world folders under the maps roots with region/*.mca but no map.xml and not already a map: the new-map import candidates (B8 ``open a local folder'' source). \\
\id{GET} & \id{/\allowbreak{}api/\allowbreak{}maps/\allowbreak{}stage-counts} & the landing cards' tallies, each counting exactly what its list shows: sketches by the layer a map holds, the other two by the stage a map stands at. \\
\addlinespace[3pt]
\multicolumn{3}{@{}l}{\textbf{\id{/shapes}} \textcolor{ink60}{\small(3)}}\\[1pt]
\midrule[0.2pt]
\id{GET} & \id{/\allowbreak{}api/\allowbreak{}shapes/\allowbreak{}catalog} & the vocabulary as cards. Unlike the browse feed this composes nothing and has no cursor: ShapeCatalog is a bounded set derived from the emitters and the tuning constants, so the whole thing is returned and the client filters in the page.\,\textcolor{ink30}{\ldots} \\
\id{GET} & \id{/\allowbreak{}api/\allowbreak{}shapes/\allowbreak{}probe} & emit one knob combination and report what the emitter said. This is the half of the catalog that a gallery cannot be: it goes through BoxFiller, so the profile check and the DockingGate run exactly as they do in composition, and a refusal is a first-class answer rather than an error.\,\textcolor{ink30}{\ldots} \\
\id{GET} & \id{/\allowbreak{}api/\allowbreak{}shapes/\allowbreak{}probe/\allowbreak{}schema} & the knob surface the panel drives: every family the emitter builds, its minimum box, which optional knobs it accepts, and whether the production menu admits it.\,\textcolor{ink30}{\ldots} \\
\addlinespace[3pt]
\multicolumn{3}{@{}l}{\textbf{\id{/compose}} \textcolor{ink60}{\small(2)}}\\[1pt]
\midrule[0.2pt]
\id{GET} & \id{/\allowbreak{}api/\allowbreak{}compose} & the browse feed. For each seed from the cursor it composes a board, reads its StructureSummary, applies the structural sieve (wool families must-include; hub / frontline forms any-of), then --- for survivors only --- evaluates, applies the score/wool sieve, and renders the SVG.\,\textcolor{ink30}{\ldots} \\
\id{POST} & \id{/\allowbreak{}api/\allowbreak{}compose/\allowbreak{}pin} & keep a browse card. Re-composes the plan from its reproducible descriptor (ComposeStages, so the structural bucket key comes for free) and saves it as a generated row (SaveGeneratedAsync, idempotent by content hash) with its structure.\,\textcolor{ink30}{\ldots} \\
\addlinespace[3pt]
\multicolumn{3}{@{}l}{\textbf{\id{/configure}} \textcolor{ink60}{\small(2)}}\\[1pt]
\midrule[0.2pt]
\id{PATCH} & \id{/\allowbreak{}api/\allowbreak{}configure/\allowbreak{}\{slug\}/\allowbreak{}exclude-island} & toggle one island's exclusion. \\
\id{GET} & \id{/\allowbreak{}api/\allowbreak{}configure/\allowbreak{}\{slug\}/\allowbreak{}state} & the current scan configuration. \\
\addlinespace[3pt]
\multicolumn{3}{@{}l}{\textbf{\id{/rules}} \textcolor{ink60}{\small(2)}}\\[1pt]
\midrule[0.2pt]
\id{GET} & \id{/\allowbreak{}api/\allowbreak{}rules} & every rule the studio can cite, with what it means and what to do about it. A refusal carries an id and one sentence about the document it was refused over. The id is stable forever and outlives the task that added it, which is what makes\,\textcolor{ink30}{\ldots} \\
\id{GET} & \id{/\allowbreak{}api/\allowbreak{}rules/\allowbreak{}terms} & every evaluator term with the band it is scoring against right now. GET /api/rules answers what a rule means; this answers the numbers. A soft term's band lives in one of two places --- stated on the term itself where it is an authored\,\textcolor{ink30}{\ldots} \\
\addlinespace[3pt]
\multicolumn{3}{@{}l}{\textbf{\id{/health}} \textcolor{ink60}{\small(1)}}\\[1pt]
\midrule[0.2pt]
\id{GET} & \id{/\allowbreak{}api/\allowbreak{}health} & Liveness probe --- confirms the API host and FastEndpoints routing are up. \\
\addlinespace[3pt]
\multicolumn{3}{@{}l}{\textbf{\id{/minecraft}} \textcolor{ink60}{\small(1)}}\\[1pt]
\midrule[0.2pt]
\id{GET} & \id{/\allowbreak{}api/\allowbreak{}minecraft/\allowbreak{}player} & ?name=\textbar{}uuid= --- resolve a Minecraft username or UUID to \{uuid, name\}. The editor uses it to turn a typed username into the canonical uuid a map.xml stores (and to resolve a stored uuid back to a name). 404 means no account is called that, which is not the same as an error\,\textcolor{ink30}{\ldots} \\
\addlinespace[3pt]
\multicolumn{3}{@{}l}{\textbf{\id{/objectives}} \textcolor{ink60}{\small(1)}}\\[1pt]
\midrule[0.2pt]
\id{GET} & \id{/\allowbreak{}api/\allowbreak{}objectives/\allowbreak{}vocabulary} & the destroyable designs (DestroyableStyles), the wool dyes (WoolColors) and what each optional field is worth unauthored (ObjectiveDefaults), served rather than copied onto the client because the client cannot reach PgmStudio.Domain.\,\textcolor{ink30}{\ldots} \\
\addlinespace[3pt]
\multicolumn{3}{@{}l}{\textbf{\id{/sketch}} \textcolor{ink60}{\small(1)}}\\[1pt]
\midrule[0.2pt]
\id{POST} & \id{/\allowbreak{}api/\allowbreak{}sketch} & originate a sketch: create a draft (geometry-less) map + its layout artifact. Returns the slug; the client navigates to /maps/\{slug\}/sketch. Body: optional \{name\} and an optional working frame \{width, depth, mode, centerX, centerZ\}.\,\textcolor{ink30}{\ldots} \\
\end{longtable}
